\section{Introduction}\subsection{Research Background and Motivation}Generative Artificial Intelligence (Gen-AI) is evolving from a niche auxiliary tool into an indispensable foundational resource across various industries. With its exponential growth in capabilities, Gen-AI is not only optimizing workflows but fundamentally reshaping the task structure of the labor market. Confronted with this radical technological transformation, higher education institutions (including comprehensive universities, vocational and technical colleges, and art academies) face unprecedented challenges: how to adjust talent cultivation models to ensure that their graduates possess enduring competitiveness and adaptability in the future workplace surrounded by AI?\paragraph{AI Usage Statement.} We disclose the usage of AI tools in this paper and in the \textbf{AI Tool Usage Report} at the end of the document.Translate the following LaTeX content:\subsection{Literature Review}Cite existing research to enhance the credibility of the model. \cite{freyosborne2017}\subsection{Problem Restatement}Part I (Professional Level: "Future Evolution" under the Influence of Gen-AI)  
The problem requires us to select one specific profession from each of the three categories: \textbf{STEM, Trade/Vocational Education, and Arts}. Based on traceable data and cited sources, we are to characterize the future changes of these professions in the context of "Gen-AI proliferation." Here, "future" should be understood as the evolution of professional forms within a defined prediction window (e.g., the next several years), including the direction and magnitude of potential reductions, stability, or expansions in job content and demand. "Current trajectory" can be interpreted as the trend along which Gen-AI penetration and usage in the industry are advancing. "Expected impact" refers to the measurable changes in job structure and practitioners' work methods brought about by this penetration. All judgments must specify the data sources used and explain the key driving factors we believe will propel professional changes.Part II (Educational Program Level: Providing Actionable Recommendations for Three Specific Institutional Programs)  
After completing the analysis at the occupational level, the problem further requires: selecting one "corresponding specific institution and specific program" for each occupation (from a university, a technical school/apprenticeship system, and an art school, respectively), and grounding the recommendations in the management decisions of these three programs. The recommendations should cover at least two core issues: first, whether the program scale should expand or contract in response to occupational changes (and how to attract more students during expansion or how to absorb potential students into other programs within the institution during contraction); second, what the program should "teach and how to handle the use of Gen-AI," such as different strategies like prohibition, limited allowance, or integration into the curriculum and training system. Here, "best support/best preparation" should be understood as: providing program decision-makers with \textbf{actionable and comparable} solutions, and using the output of mathematical models to furnish an evidence chain for these solutions, rather than remaining at the level of general principles.Part III (Evaluation Scope: Beyond "Employability" and Proposed Extrapolation Boundaries)  
The problem statement notes that "graduate employability" may not be the sole criterion for evaluating policy success; thus, teams are required to propose and justify additional evaluation dimensions, and explain how conclusions and recommendations would change when these dimensions are incorporated. To avoid conceptual ambiguity, "employability" can be understood as graduates' ability to obtain and maintain positions in a job market where Gen-AI is widely prevalent; while "other factors" may include, but are not limited to: the resource consumption costs (energy/water usage) of teaching and using Gen-AI, as well as attribution and integrity risks in academic and creative contexts. Finally, if a team believes that recommendations for a particular institution/program can be generalized to other institutions or programs, they must explicitly state the conditions and boundaries for such extrapolation, clarifying under which circumstances the generalization holds and under which it does not.
